\begin{abstract}

We introduce PULSUS, a direct frequency-domain framework for
multidimensional spectroscopy of open quantum systems beyond the
impulsive rotating-wave approximation.
The method represents coherence-frequency intervals through
Liouvillian resolvents and incorporates finite Gaussian pulses through
operator-valued functions of the field-free Liouvillian.
This construction permits spectra to be evaluated directly on
Cartesian grids or at arbitrary paired spectral coordinates while
retaining molecular evolution during the pulse envelopes and, when
required, the complete light--matter interaction.

The finite-pulse implementation is validated against an analytical
Gaussian-pulse benchmark, with a relative discrepancy of
$7.33\times10^{-6}$ at a pulse separation of $T=2\pi\sigma$.
Arbitrary-coordinate and Cartesian-grid evaluations agree to numerical
precision.
For a coupled-dimer benchmark, targeted spectral sampling reduces the
number of evaluated coordinates required for approximately $10\%$
reconstruction error by a factor of about $6.2$, with a corresponding
wall-clock reduction of about $1.8$.

A controlled hierarchy of four models is then used to separate errors
arising from the rotating-wave approximation, molecular evolution
during finite pulses, and finite spectral bandwidth.
The resulting regime map reveals a nonmonotonic dependence of the total
impulsive-RWA error on pulse duration, with the highest accuracy
occurring in an intermediate regime.
For ultrashort pulses, the dominant nonrephasing beyond-RWA
correction is traced to a coupling-enabled molecular-interaction
sequence $\mathbf{m}=(+,+,-)$.
PULSUS therefore provides both a direct computational route to
finite-pulse multidimensional spectra and a framework for diagnosing
the approximations that control them.

\end{abstract}